\section*{Chapter 10 --- The driver}
\addcontentsline{toc}{section}{Chapter 10 --- The driver}

\solution{ex:10:1}{Break the code on purpose}
The exercise asks for code changes; the outcomes are described here.
\ans{a} The trigger is written before the data registers, so the controller reads the previous
frame's contents. On the bus you see a node running exactly one frame behind --- it always sends
what you asked for \emph{last} time.
\ans{b} Data bytes 4--7 end up in \code{TX_DATA_LO}. A frame with \code{dlc = 2} and data
\code{AA BB} is sent as two zero bytes; the receiver gets the right length and the wrong content.
\ans{c} \code{AA BB CC DD} arrives as \code{DD CC BB AA}. The frame arrives, with the right length
--- just back to front.
\ans{d} The error bit never goes low, so \code{hasError()} returns \code{true} forever after the
first error and all error handling above stops working.
\ans{e} \code{STATUS} bit 1 stays set. Every \code{receive()} returns \code{true} with the same
frame, in an endless loop, and no new frame can be received.
\ans{f} Nothing shows in the output, and that is the point. \code{PrintTransport} shows which
bytes went out, and they are correct --- the fault is in what the driver \emph{hands back} to the
caller, that is, in \code{frame.data} above \code{dlc}. It is caught by a test case delivering a
frame with \code{dlc = 2} after one with \code{dlc = 8} and checking that the six topmost bytes
are zeroed, which requires a double modelling the hardware's behaviour --- that is,
\code{BankTransport} in \kapref{ch:testning}, not the scripted one.
\ans{g} Writes happen in the middle of an ongoing transmission and overwrite data the controller
is using. The outcome depends on the timing, which makes the fault sporadic.

\solution{ex:10:2}{Count transactions}
\ans{a} Six: one \code{STATUS} read and five writes.
\ans{b} One --- the \code{STATUS} read. After that it returns \code{false}.
\ans{c} \textbf{Zero.} The validation happens before everything else, so an invalid frame costs
not a single transaction.
\ans{d} One.
\ans{e} Six: \code{STATUS}, \code{RX_ID}, \code{RX_DLC}, two data registers and \code{RX_ACK}.

Question c is the point of the order: a caller that happens to send invalid frames in a loop does
not load the bus at all.

\solution{ex:10:3}{The infinite loop}
The exercise asks for code; the outcomes are described here.
\ans{a} The loop spins until the hardware becomes free, and then ends normally. That is the
intended behaviour.
\ans{b} The loop spins forever. \code{isValid()} returns \code{false} however long you wait, and
nothing in the system is going to change that.
\ans{c} Validate before the loop and return immediately on an invalid frame, and count attempts in
the loop so that even hardware that never becomes free gives up. See \secref{sec:waiting}.

\solution{ex:10:4}{Design questions}
\ans{a} An enumeration as the return value instead of \code{bool}, for example: \code{Accepted},
\code{Busy}, \code{Invalid}. The cost for \code{Stub} is small --- it would return
\code{Accepted} or \code{Invalid} --- but the cost for \emph{the interface} is that it grows, and
every caller has to handle more cases. The trade-off is that a \code{bool} is simpler to implement
and harder to use correctly.
\ans{b} A test case delivering a frame with \code{dlc = 8} and then one with \code{dlc = 2}.
Without masking, the second frame still carries six bytes from the first, and the application
cannot tell them from real data. On real hardware it also takes no more than the hardware group
changing when the data accumulator is cleared, or you borrowing a board built by another group.
\ans{c} The driver observes that \code{STATUS} bit 0 goes low at the trigger and comes back
shortly afterwards --- exactly as on a successful transmission --- and that the error flag is set.
It \emph{cannot} observe that it was an arbitration that was lost rather than a stuff error, a bad
CRC or a received remote frame, and it cannot see which node won or which frame went through
instead.
\ans{d} Polling more often, that is, doing less work between the \code{receive()} calls. The limit
you cannot get past is that a \code{receive()} takes around 270~µs while a frame takes 130~µs: a
transmitter running continuously always has time to send more frames than you have time to fetch.
